\section{Responsabilidades y colaboración}

\begin{longtable}{p{0.25\textwidth}p{0.67\textwidth}}
\toprule
\textbf{Responsabilidad} & \textbf{Funciones en este manual y en el trabajo diario} \\
\midrule
\endfirsthead
\toprule
\textbf{Responsabilidad} & \textbf{Funciones en este manual y en el trabajo diario} \\
\midrule
\endhead
Julián Carrión Tovar\\Responsable de accesibilidad & Mantener y ampliar las pautas de accesibilidad; ayudar a revisar colores, textos, interacción y alternativas; facilitar ejercicios de revisión y registrar problemas de accesibilidad. \\
\addlinespace
Yeray Rodríguez Navas\\Responsable de calidad & Mantener los criterios de calidad; promover revisiones y pruebas; ayudar a clasificar incidencias y comprobar que las correcciones cumplen los criterios acordados. \\
\addlinespace
Resto del equipo & Aplicar las pautas desde el diseño hasta la entrega, revisar cambios de compañeros, aportar pruebas y avisar cuando una decisión no esté clara o no pueda verificarse. \\
\addlinespace
Equipo completo & Acordar prioridades y criterios; validar las decisiones que afecten a las personas usuarias; registrar pendientes y actualizar este manual cuando cambie el alcance. \\
\bottomrule
\end{longtable}

\subsection{Cómo colaborar ante un desacuerdo}
Describir el problema observable, a quién puede afectar y qué criterio o necesidad está implicado. Proponer una alternativa y una forma de comprobarla. Si el equipo no puede resolverlo con la información disponible, registrar la decisión como pendiente y consultar a las personas responsables o al centro; no sustituir una necesidad no confirmada por una suposición.